\documentclass[12pt]{article}
\usepackage{graphicx} % Required for inserting images
\usepackage{amsmath}
\usepackage{amssymb}
\usepackage{indentfirst} % to get indentation on the first paragram of a section
% Set page size and margins
% Replace `letterpaper' with `a4paper' for UK/EU standard size
\usepackage[letterpaper,top=2cm,bottom=2cm,left=3cm,right=3cm,marginparwidth=1.75cm]{geometry}

\title{Combining Generalized Dubins Path Problem with visiting a circular boundary? On optimal Dubins paths for a capture the flag game? Minimal time flag capture in a capture the flag game?}
\author{Émile Marcotte}
% \date{December 2025}
\date{}

\begin{document}

% \maketitle


%%%%%%%%%%%%%%%%%%%%%%%%%%%%%%%%%%%%%%%%%%%%%%%
% \section{Potential title?}

% Combining a generalized Dubins path problem with visiting a circular boundary? On optimal Dubins paths for a capture the flag game? Minimal time flag capture in a capture the flag game?
%%%%%%%%%%%%%%%%%%%%%%%%%%%%%%%%%%%%%%%%%%%%%%%

% \section{Task at hand}

% We want to solve the following problem:

% Find the minimum time trajectory for a Dubins vehicle starting at some
% point and heading, that goes to a point on a (stationary) circle and
% ends up on the boundary of a polytope.

% Understanding those papers should enable you to answer the problem above.

% It would be any polytopic shape in the most general form, for simplicity
% we can actually start from a circle. In the context of the CTF game it is
% the robots own half.

% It is important to write the optimal control problem first. You can do it by mirroring the way people formulated the problem in those papers. Clearly define and state the initial conditions and the terminal conditions as well as the kinematics and the cost function.


\section{Pose the optimal control problem}

\subsubsection*{Problem statement}

Find the minimum time trajectory for a Dubins vehicle starting at some point and heading that goes to a point on a stationary circle and ends up on the boundary of a polytope.


\subsubsection*{Initial conditions}

The vehicle starts at a position $(x_0,y_0) \in \mathbb{R}^2$, with a heading $\theta_0 \in (\!-\!\pi,\pi]$ and a constant speed $v_0$.

We have a circular region $\mathcal{C}$ of radius $r_c$
\begin{align}
    &\mathcal{C} = \{ (x,y) \in \mathbb{R}^2 | (x-x_c)^2 + (y-y_c)^2 = r_c^2 \},\\
    &r_c \in \mathbb{R}^+, \nonumber\\
    &(x_c,y_c) \in  \mathbb{R}^2. \nonumber
\end{align}
% \begin{align}
%     &\mathcal{C} = \{ (x,y,x_c,y_c) \in \mathbb{R}^4, r_c \in \mathbb{R}^+ | (x-x_c)^2 + (y-y_c)^2 = r_c^2\}
% \end{align}

% IMAN SAID TO START WITH A CIRCLE FOR SIMPLICITY FOR THE TARGET HOME POLYTOPE
We also have a terminal area $\mathcal{X}$ that does not overlap with $\mathcal{C}$ and can be of any shape. For the sake of this example, it will be defined as a circle of radius $r_x$, noted as
\begin{align}
    &\mathcal{X} = \{ (x,y) \in \mathbb{R}^2 | (x-x_x)^2 + (y-y_x)^2 = r_x^2 \},\\
    &r_x \in \mathbb{R}^+, \nonumber\\
    &(x_x,y_x) \in  \mathbb{R}^2, \nonumber\\
    &\mathcal{C}\cap \mathcal{X} = \emptyset. \nonumber
\end{align}


\subsubsection*{Kinematics}

The kinematics of the Dubins vehicle are
\begin{align}
    &\dot{x} = v_0\cos\theta,\\
    &\dot{y} = v_0\sin\theta,\\
    &\dot{\theta} = u,\\
    &v_0\in \mathbb{R}^+,\\
    &\theta \in (\!-\!\pi,\pi],\\
    &u \in \left[-\frac{v_0}{\rho},\frac{v_0}{\rho}\right],\\
    &\rho \leq r_c,
\end{align}
where $\rho$ is the minimal turning radius of the Dubins vehicle (justify its constraint later?).

\subsubsection*{Terminal conditions}

The vehicle has to go from its initial point $(x_0,y_0)$ to a point $(x_g,y_g)$ located on the perimeter of $\mathcal{C}$, then end up at a point $(x_f,y_f)$ located on the boundary of $\mathcal{X}$ with a heading $\theta_f$ pointing inside $\mathcal{C}$.

% We need to find the terminal position $(x_T,y_T)\in \mathbb{X}$. In the context of a capture the flag game, $\mathbb{X}$ would be the boundary of the corresponding player's field, which typically forms a rectangle. If we choose to represent the field as shown in figure ADD FIGURE OF CTF FIELD, the left and right fields could be represented as follows
% \begin{align}
%     \mathbb{X}_L &= \{ (x,y)\in\mathbb{R} | -X_x \leq x \leq -\varepsilon, \;\; -X_y \leq y \leq X_y \} \\
%     \mathbb{X}_R &= \{ (x,y)\in\mathbb{R} | \; \varepsilon \leq x \leq X_x, \;\; -X_y \leq y \leq X_y \}
% \end{align}


\subsubsection*{Cost function}

The cost function is
\begin{align}
    \underset{u(t)}{\min}\int_0^T 1dt
\end{align}
where $T$ is the time it takes to go from $(x_0,y_0)$ to $(x_f,y_f)$ while transiting through $(x_g,y_g)$.

Although the objective is to find the fastest path, since the vehicle travels at a constant speed and is assumed to be in a homogeneous environment (friction and drag are constant, and no external perturbations), finding the fastest path is equivalent to finding the shortest path.
% $\square$




\section{Figuring out how to do it}

At the point $(x_g,y_g)$, to make sure the vehicle either touches the boundary of $\mathcal{C}$ or goes fully through it (whichever gives the shortest path, depends on initial conditions), the following state equality needs to be satisfied
\begin{align}
    S_g(x,y,t) = (x-x_c)^2 + (y-y_c)^2 - r_c^2 = 0.
\end{align}

At the final destination, the following state equality needs to be satisfied
\begin{align}
    S_f(x,y,t) = (x-x_c)^2 + (y-y_c)^2 - r_c^2 = 0
\end{align}
to make sure the vehicle ends up on the boundary of the terminal area $\mathcal{X}$.

At $\mathcal{C}$, the angle $\theta_g$ needs to satisfy the following contraint
$$(INSERT.CONSTRAINT)$$
(the total path length will also influence this angle's value)

The final angle $\theta_f$ needs to satisfy the following constraint
$$(INSERT .CONSTRAINT)$$
angle needs to be pointing inside $\mathcal{X}$, needs to be formally expressed.

Create the Hamiltonian for the stated problem
$$(INSERT. HAMILTONIAN).$$

From this, find the optimal control $u^*$ that minimizes the Hamiltonian. This could be a long step, because we need to take into account every constraint.



\section{Random notes}

The terminal position $(x_f,y_f)$ needs to be calculated such that this point of arrival gives the shortest distance. The "transition point" $(x_g,y_g)$ also needs to be found.

We do not have to enter the area of $\mathcal{C}$, but we do need to touch its boundary at least at one point. With that being said, it is also fine to completely go inside $\mathcal{C}$ if it allows for a quicker path to $\mathcal{X}$.

For the original Dubins algorithm: \cite{Dubins57}

To find a path that tangentially touches the circular flag area: \cite{Jha20}.

% For a DTSPN algorithm in which it is allowed for the path to intersect with the circle: \cite{Chen2021}.

To go through a sequence of points: \cite{Manyam2017}.

To catch a moving enemy: \cite{Meyer2015}.



% \bibliographystyle{alpha}
\bibliographystyle{apalike}
\bibliography{bibl}


\end{document}
